\chapter{Il linguaggio Kotlin}

\section{Variabili e Tipi Base}
In Kotlin \textbf{tutto è un oggetto} (non esistono tipi primitivi come in Java). Il compilatore non effettua conversioni implicite (es. non puoi assegnare un \texttt{Int} a un \texttt{Long} senza \texttt{.toLong()}).

\begin{itemize}
    \item \textbf{\texttt{val}:} Immutabile (sola lettura).
    \item \textbf{\texttt{var}:} Mutabile.
\end{itemize}

\section{Null Safety}
Per evitare il crash dell'app (\texttt{NullPointerException}), Kotlin introduce tipi nullable.

\begin{center}
    \begin{tikzpicture}[
        node distance=1cm and 1.5cm,
        box/.style={rectangle, draw, rounded corners, fill=white, text width=3.5cm, align=center, minimum height=1.2cm, drop shadow},
        arrow/.style={-Stealth, thick}
    ]

% Nodes
        \node[box, fill=blue!10] (nullable) {\textbf{String?}\\ (Può essere null)};
        \node[box, below=1.5cm of nullable] (safecall) {\textbf{Safe Call (?.)}\\Ritorna null se la variabile è null};
        \node[box, left=of safecall] (elvis) {\textbf{Elvis (?:)}\\Valore di default};
        \node[box, right=of safecall] (bang) {\textbf{Assertion (!!)}\\Forza: rischio crash!};

% Arrows
        \draw[arrow] (nullable) -- (safecall);
        \draw[arrow] (nullable) -| (elvis);
        \draw[arrow] (nullable) -| (bang);
    \end{tikzpicture}
\end{center}

\section{Classi e Costruttori}
L'ordine di esecuzione è un tipico trabocchetto d'esame:
\begin{enumerate}
    \item Inizializzazione delle proprietà della classe.
    \item Blocchi \textbf{\texttt{init}} (nell'ordine in cui appaiono).
    \item \textbf{Costruttore secondario} (se invocato).
\end{enumerate}

\begin{tcolorbox}[colback=red!5!white,colframe=red!75!black,title=Regola d'oro]
    Il costruttore secondario \textbf{deve sempre delegare} al primario usando \texttt{this()}.
\end{tcolorbox}

\section{Interfacce vs Classi Astratte}
\begin{table}[h]
    \centering
    \begin{tabularx}{\textwidth}{|l|X|X|}
\hline
\rowcolor{gray!20} \textbf{Caratteristica} & \textbf{Interfaccia} & \textbf{Classe Astratta} \\ \hline
\textbf{Stato} & No (no variabili istanza) & Sì (variabili istanza) \\ \hline
\textbf{Ereditarietà} & Molteplice & Singola \\ \hline
\textbf{Metodi} & Astratti o con default & Astratti o concreti \\ \hline
    \end{tabularx}
\end{table}

\section{Functional Programming}
Kotlin permette di trattare le funzioni come variabili.
\begin{itemize}
    \item \textbf{Lambda:} Funzioni anonime \texttt{\{ a, b -> a + b \}}.
    \item \textbf{Uso di \texttt{it}:} Se la lambda ha un solo parametro, puoi omettere il nome e usare la keyword \texttt{it}.
\end{itemize}

\begin{lstlisting}[caption={Esempio di Lambda e Filter}]
val numeri = listOf(3, 8, 12, 5)
// Versione estesa
val maggiori7 = numeri.filter { n -> n > 7 }
// Versione compatta con 'it'
val maggiori7Compatta = numeri.filter { it > 7 }
\end{lstlisting}